\subsubsection{Fourier Transform of Convolutions}\label{subsec:fourier-convolutions}

Another useful property of the Fourier transform is that it turns convolutions into products and vice versa. We shall first consider test functions. It is easy to see that the product and convolution of functions in $\mathcal{S}(\mathbb{R}^m)$ are again in $\mathcal{S}(\mathbb{R}^m)$. Their behavior under Fourier transform is given by the next lemma.

\begin{Lemma}
\label{lem:fourier-convolution-test-functions}
Let $\phi,\psi\in\mathcal{S}(\mathbb{R}^m)$. Then
\begin{equation}\label{eq:fourier-convolution-product-test-functions}
\mathcal{F}[\phi*\psi]=(2\pi)^{m/2}\mathcal{F}[\phi]\mathcal{F}[\psi],
\end{equation}

\begin{equation}\label{eq:fourier-product-convolution-test-functions}
\mathcal{F}[\phi\psi]=(2\pi)^{-m/2}\mathcal{F}[\phi]*\mathcal{F}[\psi].
\end{equation}

\end{Lemma}

\begin{Proof}
We have
\begin{equation}\label{eq:fourier-convolution-test-functions-proof}
\begin{aligned}
\mathcal{F}[\phi*\psi](\bm{\xi})
&=(2\pi)^{-m/2}\int_{\mathbb{R}^m}e^{-i\bm{\xi}\cdot\bm{x}}\int_{\mathbb{R}^m}\phi(\bm{x}-\bm{y})\psi(\bm{y})\,\mathrm{d}\bm{y}\,\mathrm{d}\bm{x}\\
&=(2\pi)^{-m/2}\int_{\mathbb{R}^m}\psi(\bm{y})\int_{\mathbb{R}^m}e^{-i\bm{\xi}\cdot\bm{x}}\phi(\bm{x}-\bm{y})\,\mathrm{d}\bm{x}\,\mathrm{d}\bm{y}\\
&=(2\pi)^{-m/2}\int_{\mathbb{R}^m}\psi(\bm{y})\int_{\mathbb{R}^m}e^{-i\bm{\xi}\cdot(\bm{z}+\bm{y})}\phi(\bm{z})\,\mathrm{d}\bm{z}\,\mathrm{d}\bm{y}\\
&=(2\pi)^{m/2}\mathcal{F}[\phi](\bm{\xi})\mathcal{F}[\psi](\bm{\xi}).
\end{aligned}
\end{equation}
This yields \eqref{eq:fourier-convolution-product-test-functions}. Applying the inverse Fourier transform, we can restate \eqref{eq:fourier-convolution-product-test-functions} as
\begin{equation}\label{eq:inverse-fourier-product-convolution}
\mathcal{F}^{-1}[\widehat\phi\widehat\psi]=(2\pi)^{-m/2}\mathcal{F}^{-1}[\widehat\phi]*\mathcal{F}^{-1}[\widehat\psi].
\end{equation}
This and the trivial identity
\begin{equation}\label{eq:inverse-fourier-forward-reflection}
\mathcal{F}^{-1}[\widehat\phi](\bm{x})=\mathcal{F}[\widehat\phi](-\bm{x})
\end{equation}
lead to \eqref{eq:fourier-product-convolution-test-functions}.
\end{Proof}

We shall now extend this result to distributions.

\begin{Theorem}
\label{thm:fourier-convolution-compact-distribution}
Let $f,g\in\mathcal{S}^{\prime}(\mathbb{R}^m)$ and let $g$ have compact support. Then $f*g\in\mathcal{S}^{\prime}(\mathbb{R}^m)$ and
\begin{equation}\label{eq:fourier-convolution-compact-distribution-product}
\mathcal{F}[f*g]=(2\pi)^{m/2}\mathcal{F}[f]\mathcal{F}[g].
\end{equation}

\end{Theorem}

We note that $\mathcal{F}[g]$ is a $C^\infty$ function with polynomial growth (see Example~\ref{ex:fourier-compact-distribution-entire} above) and therefore the right-hand side of \eqref{eq:fourier-convolution-compact-distribution-product} is well defined. A similar result can be established for tempered distributions on $\mathbb{R}$ whose supports are bounded from the same side; see Problem~\ref{ex:one-sided-tempered-fourier-product}.

\begin{Proof}
By definition, we have
\begin{equation}\label{eq:convolution-nested-pairing-proof}
(f*g,\phi)=(f(\bm{x}),(g(\bm{y}),\phi(\bm{x}+\bm{y}))).
\end{equation}
Since $g$ has compact support, it is of finite order, i.e., with $Q$ denoting any compact set containing the support of $g$ in its interior, there is $n\in\mathbb{N}$ and $C>0$ such that
\begin{equation}\label{eq:compact-distribution-parameter-pairing-bound}
|(g(\bm{y}),\phi(\bm{x}+\bm{y}))|\leq C\max_{\bm{y}\in Q}\sum_{|\alpha|\leq n}|D^\alpha\phi(\bm{x}+\bm{y})|.
\end{equation}
It is easy to see that, for every $k$ and $\alpha$, $|\bm{x}|^k|D^\alpha\phi(\bm{x}+\bm{y})|$ is bounded uniformly for $\bm{x}\in\mathbb{R}^m$ and $\bm{y}\in Q$; hence \eqref{eq:compact-distribution-parameter-pairing-bound} leads to a uniform bound for $|\bm{x}|^k|(g(\bm{y}),\phi(\bm{x}+\bm{y}))|$. Also, we can replace $\phi$ in \eqref{eq:compact-distribution-parameter-pairing-bound} by any of its derivatives. We conclude that the mapping $\phi\mapsto(g(\bm{y}),\phi(\bm{x}+\bm{y}))$ is continuous from $\mathcal{S}(\mathbb{R}^m)$ into itself. Hence \eqref{eq:convolution-nested-pairing-proof} is well defined and represents a continuous linear functional on $\mathcal{S}(\mathbb{R}^m)$.

From the definition of the Fourier transform, we now find
\begin{equation}\label{eq:fourier-convolution-compact-distribution-proof}
\begin{aligned}
&(\mathcal{F}[f*g](\bm{\xi}),\phi(\bm{\xi}))=(f*g(\bm{x}),\mathcal{F}^{-1}[\phi](\bm{x}))\\
&\quad=(2\pi)^{-m/2}\left(f(\bm{x}),\left(g(\bm{y}),\int_{\mathbb{R}^m}e^{i\bm{\xi}\cdot(\bm{x}+\bm{y})}\phi(\bm{\xi})\,\mathrm{d}\bm{\xi}\right)\right)\\
&\quad=\left(f(\bm{x}),\int_{\mathbb{R}^m}e^{i\bm{\xi}\cdot\bm{x}}\phi(\bm{\xi})\overline{\mathcal{F}[g](\bm{\xi})}\,\mathrm{d}\bm{\xi}\right)\\
&\quad=(2\pi)^{m/2}(f,\mathcal{F}^{-1}[\phi\overline{\mathcal{F}[g]}])\\
&\quad=(2\pi)^{m/2}(\mathcal{F}[f],\phi\overline{\mathcal{F}[g]})\\
&\quad=(2\pi)^{m/2}(\mathcal{F}[f]\mathcal{F}[g],\phi).
\end{aligned}
\end{equation}
This gives us \eqref{eq:fourier-convolution-compact-distribution-product} and completes the proof.
\end{Proof}
